\subsection{The universal obstruction for two cross arrows}

\begin{lemma}[Cancelling a trivially acting vertex in degree two]
\label{C:trivial-biset-side-H2}
Let $S$ be a nonempty finite $(N,M)$-biset, with trivial left
$N$-action.  Collapsing $N$ to the trivial monoid gives an actual
functor from its two-object collage to the ordinary right-act cone
$\mathcal C(M,S)$.  This functor induces an isomorphism on rational
degree-two cohomology.  The opposite statement holds when the right
$M$-action is trivial.
\end{lemma}
\begin{proof}
Write $E_R$ for the right-action bar of $M$ on $S$.  The edge
double bar of the collage is literally $E_R\times BN$, and its
maps to the two vertex spaces are the action-bar projection to
$BM$ and the second projection to $BN$.
Every finite monoid has zero first rational cohomology.  K\"unneth
therefore gives
\[
 \begin{aligned}
 H^1(E_R\times BN;\mathbb Q)&=H^1(E_R;\mathbb Q),\\
 H^2(E_R\times BN;\mathbb Q)&=H^2(E_R;\mathbb Q)
 \oplus H^0(E_R;\mathbb Q)\otimes H^2(BN;\mathbb Q).
 \end{aligned}
\]
In the second direct summand, pullback from $BN$ is the diagonal
map $y\mapsto 1\otimes y$, which is injective because $E_R$ is
nonempty.  Pullback from $BM$ lands in the first summand.
Consequently the kernel of the map from the two vertex
degree-two groups to the edge group is exactly
\[
 \ker\bigl(H^2(BM;\mathbb Q)\longrightarrow
                 H^2(E_R;\mathbb Q)\bigr)\oplus0.
\]
The Mayer--Vietoris sequence for the double mapping cylinder
places its degree-two cohomology in a short exact sequence with
left term $H^1(E_R;\mathbb Q)$ and right term this kernel.
The ordinary act-cone sequence has the same two end terms.
Naturality under the stated collapse functor gives the identity
on both, and hence an isomorphism in the middle.  This argument
allows disconnected $E_R$ and does not assume that either vertex
has zero second cohomology.
\end{proof}

\begin{theorem}[Two cross arrows force zero degree-two products]
\label{C:two-cross-arrows-universal}
Let $M,N$ be arbitrary finite monoids and let $S$ be a
two-element $(N,M)$-biset.  All products of rational degree-two
classes on its two-object collage vanish.  In particular, no
such collage realizes $\mathbb CP^2$.
\end{theorem}
\begin{proof}
Consider the image of the left $N$-action in the four-element
transformation monoid of $S$.

If both constant transformations occur, any commuting right
transformation fixes both points of $S$, so the right $M$-action
is trivial.  Lemma~\ref{C:trivial-biset-side-H2}, applied with
the sides reversed, identifies the collage's rational $H^2$
with that of the ordinary two-state left-act cone over $N$.
Every positive-degree product of the latter is zero by
Lemma~\ref{C:finite-line-act-cups}, applied to the opposite
monoid.  The comparison is an actual functor, so it preserves
cup products and proves the assertion in this case.

If exactly one constant transformation occurs, no swap occurs:
composing the swap with that constant would produce the other
constant.  The constant value is therefore fixed by the left
action.  Commutation with this constant also makes the same
point fixed by the right action.  Its fixed-state bar gives a
section $BM\times BN\to E_S$ of the edge projection.
Thus the direct sum of the two vertex degree-two groups
injects into $H^2(E_S;\mathbb Q)$.

If no constant transformation occurs and the left action is
trivial, Lemma~\ref{C:trivial-biset-side-H2} reduces the assertion
to the ordinary two-state right-act cone over $M$, and again
Lemma~\ref{C:finite-line-act-cups} applies.

Finally suppose the left image contains the swap and no constant.
Every commuting right transformation is then a permutation too.
The edge double bar is the two-sheeted covering of $BM\times BN$
associated to the commuting permutation actions.  Rational
transfer makes its pullback injective.  In particular the direct
sum of the vertex degree-two groups injects into the edge group.

In the two cases with an injective vertex-to-edge map,
Mayer--Vietoris says that every collage degree-two class restricts
to zero on both vertex spaces.  It therefore lifts to relative
cohomology.  Collapsing the two vertex spaces in the double
mapping cylinder gives the suspension $\Sigma(E_{S,+})$.
All positive-degree products in its reduced cohomology vanish;
their images in absolute cohomology give the required vanishing.
These cases exhaust all possible left transformation images.
\end{proof}

\begin{remark}\label{C:two-cross-strengthens-earlier-gates}
Theorem~\ref{C:two-cross-arrows-universal} removes the hypothesis
on a vertex degree-two group from
Lemma~\ref{C:two-state-biset-H2zero}.  In particular it excludes
the three two-cross-arrow profiles $(9,5,2)$, $(8,6,2)$ and
$(7,7,2)$ in the sixteen-arrow scope without any further census.
Earlier enumerations of two-state bisets remain independent
checks of the stronger statement.
\end{remark}
